%% ══════════════════════════════════════════════════════════════
%%  CHAPTER 17 — PUBLIC REST API
%% ══════════════════════════════════════════════════════════════
\chapter{Public REST API}
\label{ch:api}

\section{Overview}
\label{sec:api:overview}

\autoscholar exposes a public REST API at \texttt{/api/v1/} providing authenticated access to the platform's paper corpus, peer reviews, reports, pipeline monitoring, search, prompts, session metadata, and platform statistics.
The API is designed for programmatic access by researchers, data scientists, and third-party applications.

\section{Authentication}
\label{sec:api:auth}

\subsection{API Key Management}

\begin{description}
  \item[Key creation:] \texttt{POST /api/v1/keys} --- Open registration (no auth required). Returns the raw key exactly once.
  \item[Key listing:] \texttt{GET /api/v1/keys} --- Authenticated users can list their masked keys.
  \item[Key revocation:] \texttt{DELETE /api/v1/keys/:id} --- Revoke a specific key.
\end{description}

\subsection{Authentication Methods}

\begin{enumerate}
  \item \textbf{Bearer token:} \texttt{Authorization: Bearer sk\_as\_...}
  \item \textbf{Query parameter:} \texttt{?api\_key=sk\_as\_...}
\end{enumerate}

Keys are stored as SHA-256 hashes in the database.
The raw key is never stored.

\subsection{Scopes}

API keys support four scopes:

\begin{table}[H]
\centering
\caption{API key scopes.}
\label{tab:api-scopes}
\begin{tabular}{ll}
\toprule
\textbf{Scope} & \textbf{Permissions} \\
\midrule
\texttt{read} & Read papers, reviews, reports, stats \\
\texttt{search} & Full-text search across resources \\
\texttt{download} & Download project ZIPs and data files \\
\texttt{write} & Create prompts and reports \\
\bottomrule
\end{tabular}
\end{table}

\subsection{Rate Limiting}

In-memory per-key rate limiting:
\begin{itemize}
  \item Default: 100 requests per hour.
  \item Configurable per key at creation time.
  \item Auto-cleanup of expired rate limit windows every 10 minutes.
  \item 429 response with \texttt{Retry-After} header when exceeded.
\end{itemize}

\section{Response Format}
\label{sec:api:format}

All responses use a standardized JSON envelope:

\begin{lstlisting}[style=typescript, caption=API response envelope.]
// Success
{
  "ok": true,
  "data": { ... },
  "meta": {
    "page": 1,
    "limit": 20,
    "total": 150,
    "totalPages": 8
  }
}

// Error
{
  "ok": false,
  "error": {
    "code": "NOT_FOUND",
    "message": "Paper not found"
  }
}
\end{lstlisting}

\section{Endpoints}
\label{sec:api:endpoints}

\subsection{Papers}

\begin{description}
  \item[\texttt{GET /api/v1/papers}] List published papers.
  \begin{itemize}
    \item Filters: \texttt{domain}, \texttt{template}, \texttt{author}, \texttt{keyword}, \texttt{q} (text search).
    \item Sorting: \texttt{newest}, \texttt{oldest}, \texttt{words}, \texttt{figures}.
    \item Pagination: \texttt{page}, \texttt{limit} (max 100).
    \item Includes generation cost data (Opus 4.6 pricing).
  \end{itemize}

  \item[\texttt{GET /api/v1/papers/:id}] Get paper details.
  \begin{itemize}
    \item Full \LaTeX{} sections, BibTeX, figure count, word count.
    \item ASID, peer review score, domain classification.
  \end{itemize}

  \item[\texttt{GET /api/v1/papers/:id/latex}] Get raw \LaTeX{} source.

  \item[\texttt{GET /api/v1/papers/:id/bibtex}] Get auto-generated BibTeX citation.

  \item[\texttt{GET /api/v1/papers/domains}] List unique domains with paper counts.

  \item[\texttt{GET /api/v1/papers/templates}] List templates with paper counts.
\end{description}

\subsection{Reviews}

\begin{description}
  \item[\texttt{GET /api/v1/reviews}] List all peer reviews.
  \begin{itemize}
    \item Filters: \texttt{recommendation}, \texttt{minScore}, \texttt{maxScore}.
    \item Sorting and pagination.
  \end{itemize}

  \item[\texttt{GET /api/v1/reviews/:paperId}] Get review for a specific paper.
  \begin{itemize}
    \item Full breakdown: overall score, 6 dimension scores, section reviews, strengths, weaknesses.
  \end{itemize}

  \item[\texttt{GET /api/v1/reviews/stats/summary}] Aggregate review statistics.
  \begin{itemize}
    \item Average scores across all dimensions.
    \item Recommendation distribution.
  \end{itemize}
\end{description}

\subsection{Reports}

\begin{description}
  \item[\texttt{GET /api/v1/reports}] List shared research reports.
  \begin{itemize}
    \item Text search and pagination.
  \end{itemize}

  \item[\texttt{GET /api/v1/reports/:shareId}] Get full report.
  \begin{itemize}
    \item Answer, sources, tool results, figures.
  \end{itemize}

  \item[\texttt{GET /api/v1/reports/:shareId/sources}] Get sources only.
\end{description}

\subsection{Search}

\begin{description}
  \item[\texttt{GET /api/v1/search}] Cross-resource full-text search.
  \begin{itemize}
    \item Searches papers, reports, and prompts.
    \item Relevance scoring (title matches weighted highest).
    \item Type filtering and pagination.
  \end{itemize}

  \item[\texttt{GET /api/v1/search/papers}] Advanced paper search.
  \begin{itemize}
    \item Additional filters: domain, author, year, word count range.
  \end{itemize}
\end{description}

\subsection{Pipeline}

\begin{description}
  \item[\texttt{GET /api/v1/pipeline/runs}] List pipeline runs.
  \begin{itemize}
    \item Status filtering: queued, researching, analyzing, drafting, validating, published, rejected, failed.
    \item Pagination.
  \end{itemize}

  \item[\texttt{GET /api/v1/pipeline/runs/:id}] Get run details.
  \begin{itemize}
    \item Governor state, Popper score, prompt, session ID.
  \end{itemize}

  \item[\texttt{GET /api/v1/pipeline/stats}] Aggregate pipeline statistics.
  \begin{itemize}
    \item Average duration, score distribution, status breakdown.
  \end{itemize}

  \item[\texttt{GET /api/v1/pipeline/active}] Currently active runs.
\end{description}

\subsection{Statistics}

\begin{description}
  \item[\texttt{GET /api/v1/stats/platform}] Platform-wide metrics.
  \begin{itemize}
    \item Papers, words, figures, \LaTeX{} lines, Python executions, tokens, cost.
  \end{itemize}

  \item[\texttt{GET /api/v1/stats/domains}] Per-domain paper counts, words, figures.

  \item[\texttt{GET /api/v1/stats/timeline}] Monthly publication trends.

  \item[\texttt{GET /api/v1/stats/analytics}] Time-series metrics.
  \begin{itemize}
    \item Configurable period and limit.
    \item Metrics: requests, users, performance, tokens, tools, errors.
  \end{itemize}
\end{description}

\subsection{Sessions}

\begin{description}
  \item[\texttt{GET /api/v1/sessions}] List anonymized research sessions.
  \begin{itemize}
    \item Sorting: newest, oldest, cost, tokens.
    \item Metadata: message count, Python executions, tools used.
    \item Content is not exposed.
  \end{itemize}

  \item[\texttt{GET /api/v1/sessions/stats}] Session aggregate statistics.
  \begin{itemize}
    \item Total tokens, cost, averages, daily activity (30 days).
  \end{itemize}
\end{description}

\subsection{Prompts}

\begin{description}
  \item[\texttt{GET /api/v1/prompts}] List public prompts.
  \begin{itemize}
    \item Search and author filtering.
  \end{itemize}

  \item[\texttt{POST /api/v1/prompts}] Create a prompt (requires \texttt{write} scope).
\end{description}

\section{Pagination}
\label{sec:api:pagination}

All list endpoints support cursor-based pagination:

\begin{itemize}
  \item \texttt{page} (default: 1) --- Page number.
  \item \texttt{limit} (default: 20, max: 100) --- Items per page.
  \item Response includes \texttt{meta.total} and \texttt{meta.totalPages}.
\end{itemize}

\section{Error Codes}
\label{sec:api:errors}

\begin{table}[H]
\centering
\caption{API error codes.}
\label{tab:api-errors}
\begin{tabular}{lll}
\toprule
\textbf{HTTP} & \textbf{Code} & \textbf{Description} \\
\midrule
400 & \texttt{BAD\_REQUEST} & Invalid parameters \\
401 & \texttt{UNAUTHORIZED} & Missing or invalid API key \\
403 & \texttt{FORBIDDEN} & Insufficient scope \\
404 & \texttt{NOT\_FOUND} & Resource not found \\
429 & \texttt{RATE\_LIMITED} & Rate limit exceeded \\
500 & \texttt{INTERNAL\_ERROR} & Server error \\
\bottomrule
\end{tabular}
\end{table}

\section{Internal API Endpoints}
\label{sec:api:internal}

Beyond the public v1 API, the server exposes internal endpoints for the frontend:

\begin{description}
  \item[\texttt{POST /api/chat}] Core streaming chat endpoint (SSE).
  \item[\texttt{POST /api/speech-to-text}] Voice transcription (ElevenLabs Scribe v2).
  \item[\texttt{POST /api/upload-data}] Data file upload (CSV, Excel, Parquet, JSON, up to 100MB).
  \item[\texttt{POST /api/share}] Create shared report.
  \item[\texttt{GET /api/share/:shareId}] Retrieve shared report.
  \item[\texttt{POST /api/generate-pdf}] Markdown to PDF (Puppeteer).
  \item[\texttt{POST /api/generate-html}] Markdown to HTML.
  \item[\texttt{POST /api/generate-docx}] Markdown to DOCX (Pandoc).
  \item[\texttt{GET /api/download/:filename}] File download from downloads directory.
  \item[\texttt{GET /api/platform/stats}] Platform statistics.
  \item[\texttt{GET /api/journal/papers}] Journal paper listing.
  \item[\texttt{GET /api/repository}] Project repository browser.
  \item[\texttt{GET /api/live-sessions}] Active session listing.
  \item[\texttt{GET /api/sessions/:sessionId/stream}] SSE reconnection for live monitoring.
\end{description}

\subsection{Admin Endpoints}

Protected by admin session authentication:

\begin{itemize}
  \item Pipeline management: trigger, abort, resume, cleanup.
  \item Scheduler configuration: enable/disable, update settings.
  \item Paper management: visibility, domain reclassification, \LaTeX{} backfill.
  \item Peer review: trigger review, view results.
  \item System: cost tracking, donations, site config.
\end{itemize}
